% Part: proof-theory
% Chapter: cut-elimination
% Section: interpolation

\documentclass[../../../include/open-logic-section]{subfiles}

\begin{document}

\olfileid[id]{pt}{cut}{itp}

\olsection{Lema Maehara dan Teorema Interpolasi Craig}

Teorema interpolasi, yang dikemukakan oleh William Craig, menyatakan bahwa
jika $!A \Entails !B$, maka terdapat !!a{sentence}~$!C$ (sang
\emph{interpolan}) sedemikian sehingga $!A \Entails !C$ dan $!C \Entails
!B$. Yang penting, interpolan~$!C$ dapat dipilih sedemikian sehingga semua
!!{constant} dan !!{predicate} yang muncul dalam~$!C$ muncul baik dalam~$!A$
maupun~$!B$. (Kita mengasumsikan bahwa tidak ada !!{function} yang terlibat.)
Tanpa pembatasan ini, $!A$ dan $!B$ secara sepele dapat berfungsi sebagai
interpolan.

Teorema interpolasi berlaku bagi logika orde pertama klasik dan pernah
disebut sebagai ``sifat penting terakhir logika orde pertama'' yang
ditemukan. Teorema ini mempunyai penerapan penting dalam teori model dan
penalaran otomatis. Motivasi dan penerapan awalnya terdapat dalam filsafat
ilmu, sebab teorema tersebut dapat digunakan untuk mengkaji gagasan-gagasan
primitif yang melaluinya suatu teori dapat atau tidak dapat diaksiomatiskan.
Teorema ini juga menyiratkan teorema keterdefinisian Beth, yang menyatakan
bahwa jika suatu teori dapat mendefinisikan sebuah konsep secara implisit,
teori itu dapat mendefinisikannya secara eksplisit.

Seperti banyak hasil dalam logika matematika, teorema interpolasi dapat
dibuktikan baik dengan metode teori model maupun teori bukti. Keunggulan
metode teori-bukti ialah bahwa metode tersebut tidak hanya menunjukkan
adanya suatu interpolan, tetapi juga menyediakan algoritme yang menghitung
interpolan dari !!a{proof} bahwa $!A \Entails !B$, misalnya !!a{proof}
dalam~\Log{G3c} dari $!A \Sequent !B$.

Misalkan $\Lang L(!A)$ adalah himpunan !!{constant}s dan !!{predicate}s
dalam~$!A$, dan $\Lang L(\Gamma) = \bigcup \Setabs{\Lang L(!A)}{!A \in
\Gamma}$. Di seluruh bagian ini, kita mengasumsikan bahwa kita berurusan
dengan !!{formula}s tanpa !!{function}s.

\begin{thm}[Teorema Interpolasi Craig]\ollabel{thm:interpolation}
Misalkan !!{language} dari~$!A$ adalah $\Lang L(!A)$ dan bahasa dari $!B$
adalah $\Lang L(!B)$. Jika $\Log{G3c} \Proves !A \Sequent !B$, maka
terdapat !!a{formula}~$!C$ dalam !!{language}~$\Lang L(!C) \subseteq
\Lang L(!A) \cap \Lang L(!B)$ sedemikian sehingga $\Log{G3c} \Proves
!A \Sequent !C$ dan $\Log{G3c} \Proves !C \Sequent !B$.
\end{thm}

Kita membuktikan suatu perumuman teorema interpolasi agar dapat memanfaatkan
struktur !!{proof}s dalam~\Log{G3c}. Untuk itu, bayangkan anteseden dan
sukseden sekuen masing-masing dipisahkan menjadi dua bagian. Kita menulis
$\maeh{\Gamma_1}{\Gamma_2} \Sequent \maeh{\Delta_1}{\Delta_2}$ untuk
menunjukkan bahwa $\Gamma_1$ dan $\Delta_1$ termasuk bagian pertama,
sedangkan $\Gamma_2$ dan $\Delta_2$ termasuk bagian kedua. Teorema
interpolasi lalu segera mengikuti dari lema berikut.

\begin{lem}[Lema Maehara]\ollabel{lem:maehara}
Jika $\Log{G3c} \Proves \maeh{\Gamma_1}{\Gamma_2} \Sequent
\maeh{\Delta_1}{\Delta_2}$, maka terdapat !!a{formula}~$!C$ dalam
!!{language}~$\Lang L(!C) \subseteq \Lang L(\Gamma_1, \Delta_1) \cap
\Lang L(\Gamma_2, \Delta_2)$ sedemikian sehingga $\Log{G3c} \Proves
\Gamma_1 \Sequent \Delta_1, !C$ dan $\Log{G3c} \Proves !C, \Gamma_2
\Sequent \Delta_2$.
\end{lem}

\begin{proof}
  Andaikan $\pi \Proves \maeh{\Gamma_1}{\Gamma_2} \Sequent
  \maeh{\Delta_1}{\Delta_2}$. Kita membuktikan klaim tersebut melalui
  induksi pada $n=\pheight{\pi}$.

  Jika $n=0$, maka $\maeh{\Gamma_1}{\Gamma_2} \Sequent
  \maeh{\Delta_1}{\Delta_2}$ merupakan aksioma, dan suatu !!{formula}
  atomik~$!A$ muncul baik dalam $\Gamma_1, \Gamma_2$ maupun dalam
  $\Delta_1, \Delta_2$. Ada empat kasus, bergantung pada apakah $!A$
  termasuk bagian pertama atau kedua di kiri, dan apakah formula itu termasuk
  bagian pertama atau kedua di kanan.
  \begin{enumerate}
    \item $!A \in \Gamma_1$ dan $!A \in \Delta_1$. Kita mencari~$!C$
    sedemikian sehingga $\Gamma_1 \Sequent \Delta_1, !C$ dan $!C, \Gamma_2
    \Sequent \Delta_2$ keduanya !!{provable}. Yang pertama sudah merupakan
    aksioma terlepas dari pilihan~$!C$, sebab $!A$ muncul di kiri maupun
    kanan. Satu-satunya pilihan yang menjamin $!C, \Gamma_2 \Sequent
    \Delta_2$ !!{provable} ialah menetapkan $!C \ident \lfalse$.
    \item $!A \in \Gamma_1$ dan $!A \in \Delta_2$. Maka $\Gamma_1
    \Sequent \Delta_1, !A$ dan $!A, \Gamma_2 \Sequent \Delta_2$ keduanya
    merupakan aksioma, sehingga kita dapat memilih $!C \ident !A$.
    \item $!A \in \Gamma_2$ dan $!A \in \Delta_1$. Maka $!A, \Gamma_1
    \Sequent \Delta_1$ dan $\Gamma_2 \Sequent \Delta_2, !A$ merupakan
    aksioma, tetapi keduanya menempatkan~$!A$ pada sisi yang salah. Jadi,
    $!A$ bukan interpolannya. Namun, dengan menggunakan $\RightR{\lnot}$
    dan $\LeftR{\lnot}$, kita dapat !!{prove} $\Gamma_1 \Sequent
    \Delta_1, \lnot!A$ dan $\lnot !A, \Gamma_2 \Sequent \Delta_2$.
    \item $!A \in \Gamma_2$ dan $!A \in \Delta_2$. Latihan.
  \end{enumerate}
  Kasus ketika $\maeh{\Gamma_1}{\Gamma_2} \Sequent
  \maeh{\Delta_1}{\Delta_2}$ merupakan aksioma karena $\lfalse \in
  \Gamma_1$ atau $\lfalse \in \Gamma_2$ diserahkan sebagai latihan.

  Jika $n>0$, $\pi$ berakhir dengan inferensi logis. Kita harus membedakan
  kasus-kasus menurut aturan yang digunakan dan bagian tempat formula
  prinsipal berada. Dalam setiap kasus, kita dapat menggunakan hipotesis
  induksi bahwa lema tersebut berlaku bagi semua !!{proof}s~$\pi_1$ dengan
  $\pheight{\pi_1} < n$, bagaimanapun sekuen akhir~$\pi_1$ dipartisi menjadi
  dua bagian. Mari kita bahas beberapa kasus tertentu.
  \begin{enumerate}
    \item $\pi$ berakhir dengan $\RightR{\lnot}$ dan formula prinsipal
    $\lnot !A$. Kita mempunyai dua kasus: bergantung pada apakah $\lnot !A$
    termasuk bagian pertama atau kedua, $\pi$ berakhir dengan salah satu
    dari
    \[
    \AxiomC{}
    \RightLabel{$\pi_1$}
    \Deduce$\maeh{!A, \Gamma_1}{\Gamma_2} \fCenter
      \maeh{\Delta_1}{\Delta_2}$
    \RightLabel{\RightR{\lnot}}
    \UnaryInf$\maeh{\Gamma_1}{\Gamma_2} \fCenter
      \maeh{\Delta_1, \lnot !A}{\Delta_2}$
    \DisplayProof
    \]
    atau
    \[
    \AxiomC{}
    \RightLabel{$\pi_1$}
    \Deduce$\maeh{\Gamma_1}{!A, \Gamma_2} \fCenter
      \maeh{\Delta_1}{\Delta_2}$
    \RightLabel{\RightR{\lnot}}
    \UnaryInf$\maeh{\Gamma_1}{\Gamma_2} \fCenter
      \maeh{\Delta_1}{\Delta_2, \lnot !A}$
    \DisplayProof
    \]
    Perhatikan bahwa jika !!{formula} prinsipal~$\lnot !A$ termasuk bagian
    pertama, kita juga menempatkan !!{formula} aktif~$!A$ dalam bagian
    pertama premis. Ini adalah pilihan yang dapat kita buat. Hipotesis
    induksi juga berlaku seandainya formula aktif kita tempatkan dalam
    bagian lain, tetapi kemudian kita tidak akan dapat menemukan interpolan
    (latihan: cobalah).

    Pertama, perhatikan keadaan pertama, ketika $\lnot !A$ termasuk bagian
    pertama. Menurut hipotesis induksi, terdapat interpolan bagi sekuen akhir
    $\pi_1$, yakni !!a{formula}~$!C_1$ sedemikian sehingga keduanya
    \begin{align*}
    !A, \Gamma_1 & \Sequent \Delta_1, !C_1 &&\text{dan} &
    !C_1, \Gamma_2 & \Sequent \Delta_2
    \intertext{dapat dibuktikan. Yang kita perlukan adalah !!a{formula}~$!C$
    sedemikian sehingga keduanya berikut ini !!{provable}:}
    \Gamma_1 & \Sequent \Delta_1, \lnot !A, !C &&\text{dan} &
    !C, \Gamma_2 & \Sequent \Delta_2.
    \end{align*}
    Jelas kita dapat mengambil $!C \ident !C_1$: hipotesis induksi telah
    memberi tahu kita bahwa sekuen di kanan !!{provable}, sedangkan sekuen
    di kiri mengikuti dari sekuen yang diperoleh melalui hipotesis induksi
    dengan~$\RightR{\lnot}$. Hipotesis induksi juga memberi tahu kita bahwa
    $\Lang L(!C_1) \subseteq \Lang L(!A, \Gamma_1, \Delta_1) \cap \Lang
    L(\Gamma_2, \Delta_2)$. Karena $\Lang L(\lnot !A) = \Lang L(!A)$ dan
    $\Lang L(!C) = \Lang L(!C_1)$, kita mempunyai $\Lang L(!C_1)
    \subseteq \Lang L(\Gamma_1, \Delta_1, \lnot !A) \cap \Lang
    L(\Gamma_2, \Delta_2)$.

    Dalam keadaan kedua, situasinya sedikit berbeda karena $\lnot !A$
    termasuk bagian kedua. Kita juga menempatkan !!{formula} aktif dalam
    premis pada bagian kedua dan menerapkan hipotesis induksi pada
    \begin{align*}
    \Gamma_1 & \Sequent \Delta_1, !C_1 &&\text{dan} &
    !C_1, !A, \Gamma_2 & \Sequent \Delta_2
    \intertext{yang dapat dibuktikan. Jika aturan \RightR{\lnot} kembali
    kita terapkan pada sekuen kedua, kita memperoleh !!{proof}s dari}
    \Gamma_1 & \Sequent \Delta_1, !C_1 &&\text{dan} &
    !C_1, \Gamma_2 & \Sequent \Delta_2, \lnot !A.
    \end{align*}
    Ini kembali tepat merupakan sekuen-sekuen yang kita inginkan agar
    !!{provable} jika $!C_1$ suatu interpolan; sekali lagi, tetapkan
    $!C \ident !C_1$. Seperti sebelumnya, $\Lang L(!C_1) \subseteq \Lang
    L(\Gamma_1, \Delta_1) \cap \Lang L(\Gamma_2, \Delta_2, \lnot !A)$.

    \item $\pi$ berakhir dengan $\RightR{\lif}$ dan !!{formula} prinsipal
    $!A \lif !B$. Kembali ada dua kasus, bergantung pada apakah $!A \lif
    !B$ termasuk bagian pertama atau kedua. !!{proof}~$\pi$ berakhir dengan
    salah satu dari
    \[
    \AxiomC{}
    \RightLabel{$\pi_1$}
    \Deduce$\maeh{!A, \Gamma_1}{\Gamma_2} \fCenter
      \maeh{\Delta_1, !B}{\Delta_2}$
    \RightLabel{\RightR{\lif}}
    \UnaryInf$\maeh{\Gamma_1}{\Gamma_2} \fCenter
      \maeh{\Delta_1, !A \lif !B}{\Delta_2}$
    \DisplayProof
    \]
    atau
    \[
    \AxiomC{}
    \RightLabel{$\pi_1$}
    \Deduce$\maeh{\Gamma_1}{!A, \Gamma_2} \fCenter
      \maeh{\Delta_1}{\Delta_2, !B}$
    \RightLabel{\RightR{\lif}}
    \UnaryInf$\maeh{\Gamma_1}{\Gamma_2} \fCenter
      \maeh{\Delta_1}{\Delta_2, !A \lif !B}$
    \DisplayProof
    \]
    Sekali lagi kita telah menempatkan !!{formula}s aktif dalam premis pada
    bagian yang sama dengan !!{formula} prinsipal dalam kesimpulan. Dalam
    keadaan pertama, kita menerapkan hipotesis induksi dan memperoleh
    !!{proof}s dari
    \begin{align*}
    !A, \Gamma_1 & \Sequent \Delta_1, !B, !C_1 &&\text{dan} & !C_1, \Gamma_2 & \Sequent \Delta_2
    \intertext{Sekuen di kiri merupakan premis yang sesuai bagi
    \RightR{\lif}. Jadi, kita mempunyai sekuen-sekuen !!{provable} berikut:}
    \Gamma_1 & \Sequent \Delta_1, !A \lif !B, !C_1 &&\text{dan} & !C_1, \Gamma_2 & \Sequent \Delta_2.
    \end{align*}
    Ini berarti $!C_1$ juga merupakan interpolan bagi sekuen akhir~$\pi$,
    sehingga sekali lagi kita dapat mengambil $!C \ident !C_1$.

    Keadaan lain ditangani dengan cara yang sama.

    \item $\pi$ berakhir dengan $\LeftR{\lif}$ dan !!{formula} prinsipal
    $!A \lif !B$. Jika $!A \lif !B$ termasuk bagian pertama,
    !!{proof}~$\pi$ berakhir dengan
    \[
    \AxiomC{}
    \RightLabel{$\pi_1$}
    \Deduce$\maeh{\Gamma_1}{\Gamma_2} \fCenter
      \maeh{\Delta_1, !A}{\Delta_2}$
    \AxiomC{}
    \RightLabel{$\pi_2$}
    \Deduce$\maeh{!B, \Gamma_1}{\Gamma_2} \fCenter
      \maeh{\Delta_1}{\Delta_2}$
    \BinaryInf$\maeh{!A \lif !B, \Gamma_1}{\Gamma_2} \fCenter
      \maeh{\Delta_1}{\Delta_2}$
    \DisplayProof
    \]
    Kini hipotesis induksi memberi kita empat sekuen !!{provable}, dua bagi
    interpolan $!C_1$ dari premis kiri dan dua bagi interpolan $!C_2$ dari
    premis kedua:
    \begin{align*}
    \Gamma_1 & \Sequent \Delta_1, !A, !C_1 &&\text{(1)} &
    !C_1, \Gamma_2 & \Sequent \Delta_2 && \text{(2)}\\
    !B, \Gamma_1 & \Sequent \Delta_1, !C_2 &&\text{(3)} &
    !C_2, \Gamma_2 & \Sequent \Delta_2 && \text{(4)}.
    \end{align*}
    Sekuen (1) dan (3) dapat menjadi premis~\LeftR{\lif} jika kita
    menerapkan pelemahan dan menambahkan $!C_2$ di kanan~(1) serta $!C_1$
    di kanan~(3):
    \[
    \AxiomC{}
    \Deduce$\Gamma_1 \fCenter \Delta_1, !A, !C_1, !C_2$
    \AxiomC{}
    \Deduce$!B, \Gamma_1 \fCenter \Delta_1, !C_1, !C_2$
    \RightLabel{\LeftR{\lif}}
    \BinaryInf$!A \lif !B, \Gamma_1 \fCenter \Delta_1, !C_1, !C_2$
    \DisplayProof
    \]
    Dengan menerapkan $\RightR{\lor}$, kita memperoleh sekuen
    \[
    !A \lif !B, \Gamma_1 \fCenter \Delta_1, !C_1 \lor !C_2.
    \]
    Ini mengisyaratkan bahwa $!C_1 \lor !C_2$ mungkin merupakan interpolan
    kita dalam kasus ini. Dan memang, kita dapat !!{prove} sekuen lain yang
    diperlukan dari sekuen (2) dan~(4) yang tersisa:
    \[
    \AxiomC{}
    \Deduce$!C_1, \Gamma_2 \fCenter \Delta_2$
    \AxiomC{}
    \Deduce$!C_2, \Gamma_2 \fCenter \Delta_2$
    \RightLabel{\LeftR{\lor}}
    \BinaryInf$!C_1 \lor !C_2, \Gamma_2 \fCenter \Delta_2$
    \DisplayProof
    \]
    Tinggal memeriksa bahwa $!C_1 \lor !C_2$ berada dalam bahasa yang benar.
    Hipotesis induksi memberi kita
    \begin{align*}
    \Lang L(!C_1) & \subseteq
    \Lang L(\Gamma_1, \Delta_1, !A) \cap \Lang L(\Gamma_2, \Delta_2) &\text{dan}\\
    \Lang L(!C_2) & \subseteq
    \Lang L(!B, \Gamma_1, \Delta_1) \cap \Lang L(\Gamma_2, \Delta_2)
    \intertext{Karena}
    \Lang L(!A) & \subseteq \Lang L(!A \lif !B) &\text{dan}\\
    \Lang L(!B) & \subseteq \Lang L(!A \lif !B),
    \intertext{kita memperoleh keduanya}
    \Lang L(!C_1) & \subseteq
    \Lang L(\Gamma_1, \Delta_1, !A \lif !B) \cap \Lang L(\Gamma_2, \Delta_2) &\text{dan}\\
    \Lang L(!C_2) & \subseteq
    \Lang L(\Gamma_1, \Delta_1, !A \lif !B) \cap \Lang L(\Gamma_2, \Delta_2)
    \intertext{dan akibatnya, $\Lang L (!C_1 \lor !C_2) = {}$}
    \Lang L(!C_1) \cup \Lang L(!C_2) & \subseteq
    \Lang L(\Gamma_1, \Delta_1, !A \lif !B) \cap \Lang L(\Gamma_2, \Delta_2),
    \end{align*}
    seperti yang diperlukan.

    Dalam kasus $!A \lif !B$ termasuk bagian kedua, situasinya berbeda,
    tetapi ditangani dengan cara serupa. Hipotesis induksi kembali memberi
    kita empat sekuen !!{provable}:
    \begin{align*}
    \Gamma_1 & \Sequent \Delta_1, !C_1 &&\text{(1)} &
    !C_1, \Gamma_2 & \Sequent \Delta_2, !A && \text{(2)}\\
    \Gamma_1 & \Sequent \Delta_1, !C_2 &&\text{(3)} &
    !C_2, !B, \Gamma_2 & \Sequent \Delta_2 && \text{(4)}.
    \end{align*}
    Kini sekuen (2) dan (4)---ditambah beberapa !!{formula} yang
    dilemahkan---menjadi premis~\LeftR{\lif}:
    \[
    \AxiomC{}
    \Deduce$!C_1, !C_2, \Gamma_2 \fCenter \Delta_2, !A$
    \AxiomC{}
    \Deduce$!B, !C_1, !C_2, \Gamma_2 \fCenter \Delta_2$
    \RightLabel{\LeftR{\lif}}
    \BinaryInf$!A \lif !B, !C_1, !C_2, \Gamma_2 \fCenter \Delta_2$
    \DisplayProof
    \]
    Sekali lagi kita ingin mengubahnya menjadi sekuen dengan satu
    !!{formula} yang dibangun dari $!C_1$ dan $!C_2$, tetapi kali ini kita
    memerlukannya dalam anteseden (karena sekarang kita menangani bagian
    kedua). Penerapan $\LeftR{\land}$ menyelesaikannya; kita harus mengambil
    $!C \ident (!C_1 \land !C_2)$ sebagai interpolan. Kebetulan sekali,
    kita dapat menggunakan sekuen (1) dan (3) yang tersisa untuk !!{prove}
    sekuen yang bersesuaian bagi bagian lain:
    \[
    \AxiomC{}
    \Deduce$\Gamma_1 \fCenter \Delta_1, !C_1$
    \AxiomC{}
    \Deduce$\Gamma_1 \fCenter \Delta_1, !C_2$
    \RightLabel{\RightR{\land}}
    \BinaryInf$\Gamma_1 \fCenter \Delta_1, !C_1 \land !C_2$
    \DisplayProof
    \]
    Seperti sebelumnya, kita mempunyai $\Lang L(!C_1 \land !C_2)
    \subseteq \Lang L(\Gamma_1, \Delta_1) \cap \Lang L(\Gamma_2,
    \Delta_2, !A \lif !B)$.

    \item $\pi$ berakhir dengan $\RightR{\lforall}$ dan !!{formula}
    prinsipal~$\lforall[x][!A(x)]$:
    \[
    \AxiomC{}
    \RightLabel{$\pi_1$}
    \Deduce$\maeh{\Gamma_1}{\Gamma_2} \fCenter
    \maeh{\Delta_1, !A(c)}{\Delta_2}$
    \RightLabel{\RightR{\lforall}}
    \UnaryInf$\maeh{\Gamma_1}{\Gamma_2} \fCenter
      \maeh{\Delta_1, \lforall[x][!A(x)]}{\Delta_2}$
    \DisplayProof
    \]
    atau
    \[
    \AxiomC{}
    \RightLabel{$\pi_1$}
    \Deduce$\maeh{\Gamma_1}{\Gamma_2} \fCenter
      \maeh{\Delta_1}{\Delta_2, !A(c)}$
    \RightLabel{\RightR{\lforall}}
    \UnaryInf$\maeh{\Gamma_1}{\Gamma_2} \fCenter
      \maeh{\Delta_1}{\Delta_2, \lforall[x][!A(x)]}$
    \DisplayProof
    \]
    Dalam keadaan pertama, hipotesis induksi memberi kita !!{proof}s dari
    $\Gamma_1 \Sequent \Delta_1, !A(c), !C_1$ dan $!C_1, \Gamma_2 \Sequent
    \Delta_2$. Kita dapat menerapkan $\RightR{\lforall}$ pada yang pertama
    untuk memperoleh $\Gamma_1 \Sequent \Delta_1,
    \lforall[x][!A(x)], !C_1$. (Syarat eigenvariabel terpenuhi, karena telah
    terpenuhi pada aturan \RightR{\lforall} di akhir~$\pi$.) Jadi, $!C_1$
    akan menjadi interpolan jika syarat $\Lang L(!C_1) \subseteq \Lang
    L(\Gamma_1, \Delta_1, \lforall[x][!A(x)]) \cap \Lang L(\Gamma_2,
    \Delta_2)$ terpenuhi. Menurut hipotesis induksi, $\Lang L(!C_1)
    \subseteq \Lang L(\Gamma_1, \Delta_1, !A(c)) \cap \Lang L(\Gamma_2,
    \Delta_2)$. Jadi, kita harus mengkhawatirkan~$c$---dapatkah $c$ muncul
    dalam~$!C_1$? Tidak. Jika $c$ muncul dalam $!C_1$, berlaku baik $c \in
    \Lang L(\Gamma_1, \Delta_1, !A(c))$ (yang memang benar) maupun $c \in
    \Lang L(\Gamma_2, \Delta_2)$. Namun, ini berarti $c$ muncul dalam
    kesimpulan inferensi $\RightR{\lforall}$ dalam~$\pi$, yang melanggar
    syarat eigenvariabel.

    Keadaan lainnya serupa.

    \item $\pi$ berakhir dengan $\RightR{\lexists}$ dan !!{formula}
    prinsipal~$\lexists[x][!A(x)]$. Kembali ada dua keadaan. Kita membahas
    keadaan ketika $\lexists[x][!A(x)]$ termasuk bagian pertama dan
    menyerahkan keadaan lain sebagai latihan.

    Dalam keadaan pertama, $\pi$ berakhir dengan
    \[
    \AxiomC{}
    \RightLabel{$\pi_1$}
    \Deduce$\maeh{\Gamma_1}{\Gamma_2} \fCenter
      \maeh{\Delta_1, \lexists[x][!A(x)], !A(t)}{\Delta_2}$
    \RightLabel{\RightR{\lexists}}
    \UnaryInf$\maeh{\Gamma_1}{\Gamma_2} \fCenter
      \maeh{\Delta_1, \lexists[x][!A(x)]}{\Delta_2}$
    \DisplayProof
    \]
    Hipotesis induksi memberi kita !!{proof}s dari $\Gamma_1 \Sequent
    \Delta_1, \lexists[x][!A(x)], !A(t), !C_1$ dan $!C_1, \Gamma_2
    \Sequent \Delta_2$. Kita dapat menerapkan $\RightR{\lexists}$ pada
    yang pertama untuk memperoleh $\Gamma_1 \Sequent \Delta_1,
    \lexists[x][!A(x)], !C_1$. Jadi, $!C_1$ kembali akan menjadi interpolan
    jika syarat $\Lang L(!C_1) \subseteq \Lang L(\Gamma_1, \Delta_1,
    \lexists[x][!A(x)]) \cap \Lang L(\Gamma_2, \Delta_2)$ terpenuhi. Di
    sini hipotesis induksi memberi kita $\Lang L(!C_1) \subseteq \Lang
    L(\Gamma_1, \Delta_1, \lexists[x][!A(x)], !A(t)) \cap \Lang
    L(\Gamma_2, \Delta_2)$. Ingatlah bahwa kita mengasumsikan tidak ada
    !!{function}s. Artinya, suku~$t$ hanya dapat berupa !!a{constant}, yakni
    $t \ident b$. Jika $b \notin \Lang L(!C_1)$ atau $b \in \Lang
    L(\Gamma_1, \Delta_1, \lexists[x][!A(x)]) \cap \Lang L(\Gamma_2,
    \Delta_2)$, tidak ada masalah: kita dapat mengambil $!C_1$ sebagai
    interpolan. Namun, bagaimana jika $b \in \Lang L(!C_1)$ dan $b \notin
    \Lang L(\Gamma_1, \Delta_1, \lexists[x][!A(x)]) \cap \Lang
    L(\Gamma_2, \Delta_2)$? Pertama-tama, dengan mengganti semua kemunculan
    $b$ dalam~$!C_1$ dengan !!a{variable}~$y$, kita dapat menulis
    $!C_1 \ident \Subst{!C_1'}{b}{y}$ dengan $!C_1'$ tidak memuat~$b$.
    Selain itu, berlaku $b \notin \Lang L(\Gamma_1, \Delta_1,
    \lexists[x][!A(x)])$ atau $b \notin \Lang L(\Gamma_2, \Delta_2)$.
    Masing-masing, kita dapat mengambil $!C \ident \lforall[y][!C_1'(y)]$
    atau $!C \ident \lexists[y][!C_1'(y)]$. Dalam kasus pertama, kita
    mempunyai:
    \[
    \AxiomC{}
    \Deduce$\Gamma_1 \fCenter
      \Delta_1, \lexists[x][!A(x)], !A(b), !C_1'(b)$
    \RightLabel{\RightR{\lexists}}
    \UnaryInf$\Gamma_1 \fCenter \Delta_1, \lexists[x][!A(x)], !C_1'(b)$
    \RightLabel{\RightR{\lforall}}
    \UnaryInf$\Gamma_1 \fCenter
      \Delta_1, \lexists[x][!A(x)], \lforall[y][!C_1'(y)]$
    \DisplayProof
    \]
    dan
    \[
    \AxiomC{}
    \Deduce$!C_1'(b), \lforall[y][!C_1'(y)], \Gamma_2 \fCenter \Delta_2$
    \RightLabel{\LeftR{\lforall}}
    \UnaryInf$\lforall[y][!C_1'(y)], \Gamma_2 \fCenter \Delta_2$
    \DisplayProof
    \]
    !!{proof}s dari sekuen teratas berasal dari hipotesis induksi (pada
    yang kedua, kita menerapkan dapat diterimanya pelemahan dengan
    $\lforall[y][!C_1'(y)]$ dalam anteseden). Syarat eigenvariabel pada
    \RightR{\lforall} dalam !!{proof} pertama terpenuhi karena $b \notin
    \Lang L(\Gamma_1, \Delta_1, \lexists[x][!A(x)])$, dan $!C_1'$ telah
    didefinisikan agar tidak memuat~$b$.

    Dalam kasus kedua, kita mempunyai (dari hipotesis induksi ditambah dapat
    diterimanya pelemahan dengan~$\lexists[y][!C_1'(y)]$):
    \[
    \AxiomC{}
    \Deduce$\Gamma_1 \fCenter
      \Delta_1, \lexists[x][!A(x)], !A(b), \lexists[y][!C_1'(y)], !C_1'(b)$
    \RightLabel{\RightR{\lexists}}
    \UnaryInf$\Gamma_1 \fCenter
      \Delta_1, \lexists[x][!A(x)], \lexists[y][!C_1'(y)], !C_1'(b)$
    \RightLabel{\RightR{\lexists}}
    \UnaryInf$\Gamma_1 \fCenter
      \Delta_1, \lexists[x][!A(x)], \lexists[y][!C_1'(y)]$
    \DisplayProof
    \]
    dan
    \[
    \AxiomC{}
    \Deduce$!C_1'(b), \Gamma_2 \fCenter \Delta_2$
    \RightLabel{\LeftR{\lexists}}
    \UnaryInf$\lexists[y][!C_1'(y)], \Gamma_2 \fCenter \Delta_2$
    \DisplayProof
    \]
    Syarat eigenvariabel pada \LeftR{\lexists} dalam !!{proof} kedua
    terpenuhi karena $b \notin \Lang L(\Gamma_2, \Delta_2)$, dan $!C_1'$
    telah didefinisikan agar tidak memuat~$b$.
  \end{enumerate}
  Kasus-kasus lain serupa; kita menyerahkannya sebagai latihan.
\end{proof}

\begin{prob}
Andaikan kita mempunyai penghubung $\lnand$ dengan aturan
\[
\Axiom$\Gamma \fCenter \Delta, !A$
\Axiom$\Gamma \fCenter \Delta, !B$
\RightLabel{\LeftR{\lnand}}
\BinaryInf$!A \lnand !B, \Gamma \fCenter \Delta$
\DisplayProof
\quad
\Axiom$!A, !B, \Gamma \fCenter \Delta$
\RightLabel{\RightR{\lnand}}
\UnaryInf$\Gamma \fCenter \Delta, !A \lnand !B$
\DisplayProof
\]
Tunjukkan bahwa Lema Maehara tetap berlaku dengan mengerjakan kasus-kasus
dalam bukti \olref{lem:maehara} ketika $\pi$ berakhir dengan salah satu
aturan ini.
\end{prob}

Setelah menetapkan Lema Maehara, kita dapat menggunakannya untuk membuktikan
Teorema Interpolasi Craig.

\begin{proof}[Bukti \olref{thm:interpolation}]
  Andaikan $!A \Sequent !B$ !!{provable}. Terapkan \olref{lem:maehara} pada
  $\maeh{!A}{\ } \Sequent \maeh{\ }{!B}$ untuk memperoleh interpolan~$!C$.
  Kita mempunyai $\Proves !A \Sequent !C$ dan $\Proves !C \Sequent !B$.
\end{proof}

Andaikan $\Gamma$ suatu teori (himpunan !!{sentence}s) dalam
!!a{language}~$\Lang L$ yang memuat suatu predikat $n$-tempat~$R$. Kita
mengatakan bahwa $\Gamma$ \emph{mendefinisikan secara implisit}~$R$ jika dan
hanya jika
\begin{align*}
\Gamma, \Gamma'
& \Entails \lforall[x_1][\dots\lforall[x_n][(R(x_1,\dots,x_n) \liff
R'(x_1, \dots, x_n))]],
\intertext{dengan $\Gamma'$ sebagai $\Gamma$ setelah semua kemunculan~$R$
diganti oleh~$R' \notin \Lang L$. Teori $\Gamma$ \emph{mendefinisikan
secara eksplisit}~$R$ jika dan hanya jika terdapat suatu~$!A(x_1, \dots,
x_n)$ yang tidak memuat~$R$ sedemikian sehingga}
\Gamma &\Entails
\lforall[x_1][\dots\lforall[x_n][(R(x_1,\dots,x_n) \liff !A(x_1,
\dots, x_n))]].
\end{align*}
Jelas bahwa jika $\Gamma$ mendefinisikan $R$ secara eksplisit, maka teori
itu mendefinisikan~$R$ secara implisit. Konversnya tidak jelas, tetapi
tetap berlaku:

\begin{lem}[Teorema Keterdefinisian Beth]
  Jika $\Gamma$ mendefinisikan~$R$ secara implisit, maka teori itu juga
  mendefinisikannya secara eksplisit.
\end{lem}

\begin{proof}
  Andaikan $\Gamma$ mendefinisikan~$R$ secara implisit. Misalkan $R'$ dan
  $\Gamma'$ seperti di atas, serta $\Lang L' = \Lang L(\Gamma') = \Lang L
  \setminus \{R\} \cup \{R'\}$. Menurut hipotesis, kita mempunyai
  \begin{align*}
    \Gamma, \Gamma' &
    \Sequent \lforall[x_1][\dots\lforall[x_n][(R(x_1,\dots,x_n)
    \liff R'(x_1, \dots, x_n))]]
  \intertext{Misalkan $c_1$, \dots,~$c_n$ adalah !!{constant}s yang tidak
  muncul dalam~$\Gamma$. Menurut lema keterbalikan, kita mempunyai}
    \Gamma, \Gamma', R(c_1,\dots,c_n) & \Sequent R'(c_1, \dots, c_n)
  \intertext{Terapkan Lema Maehara pada}
    \maeh{\Gamma, R(c_1,\dots,c_n)}{\Gamma'} & \Sequent
      \maeh{\ }{R'(c_1, \dots, c_n)}
  \intertext{untuk menemukan interpolan~$!A(c_1, \dots, c_n)$. Dengan
  mengganti konstanta-konstanta baru itu secara seragam oleh variabel bebas
  yang bersesuaian, kita memperoleh formula dasarnya dan sekuen-sekuen}
    \Gamma, R(c_1,\dots,c_n) & \Sequent !A(c_1, \dots, c_n) \\
    !A(c_1, \dots, c_n), \Gamma' & \Sequent R'(c_1, \dots, c_n)
  \intertext{sebagai sekuen-sekuen yang mempunyai !!{proof}s. Pembatasan
  bahasa dalam Lema Maehara memberikan $\Lang L(!A) \subseteq
  \Lang L(\Gamma) \setminus \{R\}$ setelah konstanta baru diganti oleh
  variabel; khususnya, $!A$ tidak memuat~$R$. Ganti $R'$ dengan~$R$ di
  mana-mana dalam !!{proof} sekuen kedua untuk memperoleh !!a{proof} dari}
    !A(c_1, \dots, c_n), \Gamma & \Sequent R(c_1, \dots, c_n).
  \end{align*}
  Terapkan \RightR{\lif}, \RightR{\land}, dan \RightR{\lforall} untuk
  memperoleh !!a{proof} dari
  \[
    \Gamma \Sequent \lforall[x_1][\dots\lforall[x_n][(R(x_1,\dots,x_n)
    \liff !A(x_1, \dots, x_n))]].
  \]
  Ini menunjukkan bahwa $!A(x_1, \dots, x_n)$ mendefinisikan $R$ secara
  eksplisit dalam~$\Gamma$.
\end{proof}

Hasil ketiga yang ekuivalen dengan interpolasi (dan, seperti interpolasi,
dapat dibuktikan dengan Lema Maehara) adalah teorema berikut: Jika dua teori
konsisten~$\Gamma_1$, $\Gamma_2$ memuat suatu teori lengkap $\Gamma$ dalam
bahasa~$\Lang L(\Gamma) \subseteq \Lang L(\Gamma_1) \cap \Lang
L(\Gamma_2)$, maka $\Gamma_1 \cup \Gamma_2$ konsisten.

Ingat bahwa teori lengkap adalah teori sedemikian sehingga untuk setiap
!!{sentence}~$!A$ dalam bahasanya berlaku $\Gamma \Proves !A$ atau $\Gamma
\Proves \lnot !A$. Karena $\Gamma_1$ dan~$\Gamma_2$ merupakan perluasan
konsisten dari~$\Gamma$, $\Gamma$ harus konsisten. Akibatnya, jika $!A$
adalah !!a{sentence} dalam bahasa $\Lang L(\Gamma_1) \cap \Lang
L(\Gamma_2)$, tidak mungkin berlaku $\Gamma_1 \Entails !A$ sekaligus
$\Gamma_2 \Entails \lnot !A$. Untuk melihatnya, andaikan ada~$!A$ semacam
itu. Jika $\Gamma_1 \Entails !A$, maka $\Gamma \Entails !A$. Jika tidak,
menurut kelengkapan~$\Gamma$, berlaku $\Gamma \Entails \lnot !A$, dan
karena $\Gamma \subseteq \Gamma_1$, berlaku $\Gamma_1 \Entails \lnot !A$,
yang membuat $\Gamma_1$ tak konsisten. Jadi $\Gamma \Entails !A$, dan
karena $\Gamma \subseteq \Gamma_2$, berlaku $\Gamma_2 \Entails !A$. Maka
$\Gamma \Entails/ \lnot !A$, sebab jika tidak, $\Gamma_2$ tak konsisten.

Ketiadaan !!a{sentence}~$!A$ dalam !!{language} gabungan~$\Lang
L(\Gamma_1) \cap \Lang L(\Gamma_2)$ dengan $\Gamma_1 \Entails !A$ dan
$\Gamma_2 \Entails \lnot !A$ adalah semua yang diperlukan bagi bukti,
yang kita berikan di sini secara teori-bukti dengan menggunakan Lema
Maehara.

\begin{thm}[Teorema Konsistensi Gabungan Robinson]
Andaikan $\Gamma_1$, $\Gamma_2$ merupakan teori konsisten dan bagi tidak
ada~$!A$ dalam !!{language} $\Lang L(\Gamma_1) \cap \Lang L(\Gamma_2)$
berlaku $\Gamma_1 \Entails !A$ dan $\Gamma_2 \Entails \lnot !A$. Maka
$\Gamma_1 \cup \Gamma_2$ konsisten.
\end{thm}

\begin{proof}
  Demi kontradiksi, andaikan $\Gamma_1 \cup \Gamma_2$ tak konsisten. Maka sekuen
  $\Gamma_1, \Gamma_2 \Sequent \ $ mempunyai !!a{proof}. Terapkan Lema
  Maehara pada $\maeh{\Gamma_1}{\Gamma_2} \Sequent \maeh{\ }{\ }$ untuk
  memperoleh !!a{sentence}~$!A$ dalam !!{language}~$\Lang L(\Gamma_1)
  \cap \Lang L(\Gamma_2)$ sedemikian sehingga $\Proves \Gamma_1 \Sequent
  !A$ dan $\Proves !A, \Gamma_2 \Sequent \quad$. Dari yang terakhir, kita
  memperoleh $\Proves \Gamma_2 \Sequent \lnot !A$. Namun, kita telah
  mengasumsikan bahwa tidak ada !!{sentence} semacam itu.
\end{proof}

Di atas kita secara diam-diam mengasumsikan bahwa teori-teori $\Gamma$,
$\Gamma_1$, dan $\Gamma_2$ berhingga. Menurut kekompakan, hasil-hasil
tersebut juga berlaku bagi teori-teori tak hingga.

\end{document}
